\section{Product ancillary sites and retained prefixes}

The length parameter in the successive construction can be increased while
preserving the local dimension and a fixed positive lower gap. The following
identities describe the added on-site penalties and the extraction of the
original prefix from a viable witness.

\begin{definition}[Linear placement of an operator]
    \label{def:lvv_placement_linear_map}
    \lean{QuantumCircuit.embedOpLinearMap}
    \leanok
    \uses{lem:qc_placed_operators}
    For a map $e$ between finite sets of sites, let $\Phi_e(X)$ be the
    operator $X$ placed by $e$, with the matrix entries of
    Lemma~\ref{lem:qc_placed_operators}. The map $\Phi_e$ is complex-linear.
\end{definition}

\begin{lemma}[Order under placement of sites]
    \label{thm:lvv_placement_order}
    \lean{QuantumCircuit.embedOp_sub,QuantumCircuit.embedOp_posSemidef,
        QuantumCircuit.embedOp_le,QuantumCircuit.siteOp_posSemidef,
        QuantumCircuit.siteOp_le}
    \leanok
    \uses{def:lvv_placement_linear_map}
    Placement respects subtraction. For injective $e$, it preserves positive
    semidefiniteness and operator inequalities. In particular, placing an
    on-site positive contraction at any site gives a positive contraction.
\end{lemma}
\begin{proof}
    \leanok
    For $X\ge0$, write $X=(\sqrt X)^\dagger\sqrt X$.
    Placement preserves products and adjoints, hence
    $\Phi_e(X)=\Phi_e(\sqrt X)^\dagger\Phi_e(\sqrt X)\ge0$.
    Apply this to the positive difference in an operator inequality.
    Single-site placement is an injective placement followed by a relabelling
    of its single local coordinate.
\end{proof}

\begin{definition}[Penalties for a product basis vector]
    \label{def:lvv_product_basis_penalty}
    \lean{QuantumCircuit.siteBasisPenalty,QuantumCircuit.productBasisPenalty}
    \leanok
    On a finite set of $q$-dimensional sites, fix a basis configuration $b$.
    Let $P_v$ be the on-site matrix diagonal in the standard basis, with
    diagonal value zero at $b_v$ and one at every other basis level. Define
    $A_b=\sum_v(P_v)_v$, with each penalty placed on its actual site.
\end{definition}

\begin{lemma}[The actual product-basis penalty]
    \label{thm:lvv_product_basis_penalty}
    \lean{QuantumCircuit.siteBasisPenalty_posSemidef,
        QuantumCircuit.siteBasisPenalty_le_one,
        QuantumCircuit.productBasisPenalty_eq_diagonal,
        QuantumCircuit.productBasisPenalty_gap,
        QuantumCircuit.productBasisPenalty_mulVec}
    \leanok
    \uses{def:lvv_product_basis_penalty}
    Each $P_v$ is a positive contraction. In the full configuration basis,
    $A_b$ is diagonal with entry $\#\{v:x_v\ne b_v\}$ at configuration $x$.
    Consequently
    \begin{align*}
        A_b|b\rangle&=0,& A_b&\ge I-|b\rangle\langle b|.
    \end{align*}
    These assertions also hold for the empty set of sites.
\end{lemma}
\begin{proof}
    \leanok
    Placing a diagonal on-site operator gives the corresponding function of
    the full configuration's coordinate. Summing gives the displayed count.
    The count is zero at $b$ and at least one at every different configuration.
\end{proof}

\begin{definition}[The padded nearest-neighbour interaction]
    \label{def:lvv_padded_interaction}
    \lean{TNLean.MPS.AreaLaw.paddedInteraction,TNLean.MPS.AreaLaw.paddingBondEquiv}
    \leanok
    Let $1\le L\le N$, fix a basis level $z$, and let $h_j$ be interactions
    on the first $L$ sites. For $0\le j<N-1$, define
    \begin{align*}
        \widehat h_j=
        \begin{cases}
            \Phi(h_j),&j<L-1,\\
            (I-|z\rangle\langle z|)_{j+1},&j\ge L-1,
        \end{cases}
    \end{align*}
    where $\Phi$ places the original sites in the initial segment.
    The bond labels split bijectively into the original bond labels and the
    appended sites; the latter correspondence is $j\mapsto j+1$.
\end{definition}

\begin{lemma}[Interval supports under prefix placement]
    \label{thm:lvv_prefix_support}
    \lean{TNLean.MPS.AreaLaw.isSupportedOn_embedOp_prefix}
    \leanok
    \uses{lem:qc_placed_operators}
    Placing an operator from the first $L$ sites into the first $L$ sites of
    a longer chain preserves each interval support $[a,b)$.
\end{lemma}
\begin{proof}
    \leanok
    The placement sends the algebra on a set of original sites into the
    algebra on its image. The image of the original interval is contained
    in the same interval of the longer chain.
\end{proof}

\begin{lemma}[Support, order and sum of the padded interaction]
    \label{thm:lvv_padded_interaction_sum}
    \lean{TNLean.MPS.AreaLaw.paddedInteraction_supported,
        TNLean.MPS.AreaLaw.paddedInteraction_nonneg_le_one,
        TNLean.MPS.AreaLaw.sum_paddedInteraction}
    \leanok
    \uses{def:lvv_padded_interaction,thm:lvv_prefix_support,
        thm:lvv_placement_order,thm:lvv_product_basis_penalty}
    If $0\le h_j\le I$ is supported on $\{j,j+1\}$, then each
    $0\le\widehat h_j\le I$ is supported on $\{j,j+1\}$ and
    \begin{align*}
        \sum_{j=0}^{N-2}\widehat h_j
          =\Phi\left(\sum_{j=0}^{L-2}h_j\right)
           +\sum_{v=L}^{N-1}(I-|z\rangle\langle z|)_v.
    \end{align*}
    The interaction is specified before choosing a ground vector or energy.
\end{lemma}
\begin{proof}
    \leanok
    Original terms retain support and order under prefix placement.
    Each later penalty acts only on the right endpoint of its assigned bond.
    The bond-label bijection separates the two sums, and linearity of
    placement combines the original terms.
\end{proof}

\begin{lemma}[Norms of prefix columns]
    \label{thm:lvv_prefix_column_norm}
    \lean{TNLean.MPS.AreaLaw.sum_norm_cutColumn_sq,
        TNLean.MPS.AreaLaw.norm_cutColumn_le}
    \leanok
    Write a chain vector as a coefficient matrix $M$ across a prefix cut,
    and let $w_b$ be its column at right configuration $b$. Then
    \begin{align*}
        \sum_b\|w_b\|^2=\|\psi\|^2,\qquad \|w_b\|\le\|\psi\|.
    \end{align*}
\end{lemma}
\begin{proof}
    \leanok
    The configuration equivalence reindexes the squared norm as the double
    sum of squared coefficient moduli. Every summand is nonnegative.
\end{proof}

\begin{definition}[Extension by a fixed right basis configuration]
    \label{def:lvv_right_basis_extension}
    \lean{TNLean.MPS.AreaLaw.rightBasisExtension}
    \leanok
    For a left vector $v$ and right basis configuration $b$, let
    $E_bv=v\otimes|b\rangle$. In prefix coordinates its only nonzero
    column is the column $v$ at $b$.
\end{definition}

\begin{lemma}[Coordinates and pairing of a basis extension]
    \label{thm:lvv_right_basis_extension}
    \lean{TNLean.MPS.AreaLaw.cutColumn_rightBasisExtension,
        TNLean.MPS.AreaLaw.inner_rightBasisExtension,
        TNLean.MPS.AreaLaw.norm_rightBasisExtension}
    \leanok
    \uses{def:lvv_right_basis_extension,thm:lvv_prefix_column_norm}
    The column of $E_bv$ at $b'$ is $v$ if $b'=b$ and zero otherwise.
    Moreover $\|E_bv\|=\|v\|$ and
    $\langle\psi,E_bv\rangle=\langle w_b,v\rangle$.
\end{lemma}
\begin{proof}
    \leanok
    Evaluate the extension on the glued left and right configuration.
    The inner-product sum has only the prescribed right column, and the
    column norm identity gives preservation of the norm.
\end{proof}

\begin{lemma}[A unit vector in the retained prefix span]
    \label{thm:lvv_retained_prefix_extraction}
    \lean{TNLean.MPS.AreaLaw.IsViable.exists_prefix_of_rightBasisExtension}
    \leanok
    \uses{thm:lvv_right_basis_extension,thm:lvv_prefix_column_norm}
    Let $\delta<1$ and suppose $S$ is $\delta$-viable for $E_bv$.
    Then there is a unit vector $w\in\operatorname{span}S$ with
    $|\langle w,v\rangle|\ge1-\delta$.
    No normalization hypothesis on $v$ is required for this implication.
\end{lemma}
\begin{proof}
    \leanok
    Choose the actual unit viable witness $\psi$ and its column $w_b$.
    The support condition gives $w_b\in\operatorname{span}S$; the pairing
    identity gives $|\langle w_b,v\rangle|\ge1-\delta>0$.
    Thus $w_b\ne0$ and $\|w_b\|\le1$.
    The vector $w=w_b/\|w_b\|$ remains in the same span and has at least
    the original overlap.
\end{proof}

\begin{lemma}[Tensor coordinates of the padded terms]
    \label{thm:lvv_padding_tensor_coordinates}
    \lean{TNLean.MPS.AreaLaw.restrict_prefix_cutEquiv_symm,
        TNLean.MPS.AreaLaw.cutOperator_embedOp_prefix,
        TNLean.MPS.AreaLaw.cutOperator_appended_penalty}
    \leanok
    \uses{thm:lvv_padded_interaction_sum,thm:lvv_product_basis_penalty}
    Across the cut after the original $L$ sites, the embedded original
    Hamiltonian is $H\otimes I$, after numbering its left sites in their
    original order. The sum of the appended-site penalties is $I\otimes A_b$,
    where $b$ is the constant configuration with basis level $z$.
\end{lemma}
\begin{proof}
    \leanok
    Each operator is supported on its stated factor. Its matrix slice at
    a fixed complementary configuration identifies that factor: restriction
    of a glued configuration gives the original left configuration, and
    the right diagonal entries count the sites differing from $b$.
\end{proof}

\begin{definition}[The padded ground vector]
    \label{def:lvv_padded_ground_vector}
    \lean{TNLean.MPS.AreaLaw.paddedGroundVector}
    \leanok
    \uses{def:lvv_right_basis_extension}
    For an original chain vector $\Gamma$, define
    $\widehat\Gamma=\Gamma\otimes|z\rangle^{\otimes(N-L)}$,
    using the prefix-coordinate identification.
\end{definition}

\begin{lemma}[A projector gap for the actual padded chain]
    \label{thm:lvv_padding_ground_gap}
    \lean{TNLean.MPS.AreaLaw.norm_paddedGroundVector,
        TNLean.MPS.AreaLaw.isGappedGroundState_paddedInteraction}
    \leanok
    \uses{def:lvv_padded_ground_vector,thm:lvv_padding_tensor_coordinates,
        thm:lvv_product_basis_penalty,thm:lvv_right_basis_extension}
    Suppose $\Gamma$ is a unit vector, $H\Gamma=E\Gamma$, and
    $H-EI\ge\Delta(I-|\Gamma\rangle\langle\Gamma|)$ for $\Delta>0$.
    Set $g=\min\{\Delta,1\}$.
    The actual padded interaction has unit ground vector
    $\widehat\Gamma$ at energy $E$, with
    \begin{align*}
        \widehat H-EI\ge g(I-|\widehat\Gamma\rangle
            \langle\widehat\Gamma|).
    \end{align*}
    This is a lower bound for the projector gap.
\end{lemma}
\begin{proof}
    \leanok
    In the coordinates of Lemma~\ref{thm:lvv_padding_tensor_coordinates},
    $\widehat H=H\otimes I+I\otimes A_b$ and
    $\widehat\Gamma=\Gamma\otimes|b\rangle$.
    Both factors have projector gap at least $g$.
    Write $P=|\Gamma\rangle\langle\Gamma|$ and
    $Q=|b\rangle\langle b|$. The defect is the sum of the positive matrices
    \begin{align*}
        &(H-EI-g(I-P))\otimes I
          +I\otimes(A_b-g(I-Q))\notag\\
        &\hspace{4em}+g(I-P)\otimes(I-Q).
    \end{align*}
    The eigen-equation follows by applying the two factors separately.
    The coordinate identifications preserve positivity and normalization.
\end{proof}

\begin{lemma}[A phase-aligned prefix MPS]
    \label{thm:lvv_padding_prefix_mps}
    \lean{TNLean.MPS.AreaLaw.IsViable.exists_prefix_mps_of_paddedGroundVector}
    \leanok
    \uses{thm:lvv_retained_prefix_extraction,thm:lvv_left_span_bond,
        def:lvv_padded_ground_vector}
    Let $\Gamma$ be a unit original-chain vector and $\delta<1$.
    Suppose a finite prefix family $S$ is $\delta$-viable for
    $\widehat\Gamma$, and every member has effective MPS bond bound
    $\max\{1,b\}$. There is a unit $w\in\operatorname{span}S$ whose
    original-coordinate vector $\psi$ satisfies
    \begin{align*}
        \|\psi-\Gamma\|^2\le2\delta,
        \qquad D(\psi)\le |S|\max\{1,b\}.
    \end{align*}
\end{lemma}
\begin{proof}
    \leanok
    Apply Lemma~\ref{thm:lvv_retained_prefix_extraction} to the fixed
    ancillary configuration. Rotate the resulting unit vector by a phase
    minimizing its distance from the original target. This rotation preserves
    its span membership. The squared distance is
    $2-2|\langle\psi,\Gamma\rangle|\le2\delta$.
    The finite-span MPS bound gives the asserted bond dimension.
\end{proof}

\begin{theorem}[Uniform polynomial arbitrary accuracy from the successive construction]
    \label{thm:lvv_successive_arbitrary_accuracy}
    \lean{TNLean.MPS.AreaLaw.lvvSuccessiveError_min_lt_one,
        TNLean.MPS.AreaLaw.exists_uniform_positive_chain_successive_accuracy}
    \leanok
    \uses{thm:lvv_uniform_successive_approximation,
        thm:lvv_padded_interaction_sum,thm:lvv_padding_ground_gap,
        thm:lvv_padding_prefix_mps}
    Fix $q\ge2$ and $\Delta>0$, and put
    $g=\min\{\Delta,1\}$ and $c=(g/169)^2<1$.
    There are natural numbers $P,d$, with $P\ge1$, depending only on
    $q,\Delta$, such that the following holds.
    Let $L,R\ge1$ be natural numbers, and let the $L$-site interactions
    $0\le h_j\le I$ be supported on $\{j,j+1\}$.
    Set $N=L+2R^2$ and form the actual padded interaction of
    Definition~\ref{def:lvv_padded_interaction}.
    Its successive construction chooses one family $S_i$ at every stage
    before any original ground energy or vector, with $S_0=\{1\}$,
    cardinality at most $P(N+1)^d$, and effective member bond bound
    $P(N+1)^d$.
    For every unit $\Delta$-gapped original ground vector $\Gamma$,
    each $S_i$ is $c/N$-viable for $\widehat\Gamma$.
    Moreover its actual retained stage-$L$ span contains a unit vector $w$
    whose original-coordinate vector $\psi$ satisfies
    \begin{align*}
        \|\psi-\Gamma\|^2\le R^{-2},
        \qquad D(\psi)\le P(N+1)^d.
    \end{align*}
    This is a product-ancilla consequence of the proved successive
    construction. It asserts no computational running time or sharper
    printed quantitative bounds.
\end{theorem}
\begin{proof}
    \leanok
    Apply Theorem~\ref{thm:lvv_uniform_successive_approximation} once at
    the fixed gap $g$. If its uniform polynomial budget is $P_0(N+1)^{d_0}$,
    take $P=P_0^2$ and $d=2d_0$.
    The actual padded chain satisfies its interaction hypotheses, and
    Lemma~\ref{thm:lvv_padding_ground_gap} supplies the required
    ground vector and gap. Use the stage-$L$ family of this same run.
    Lemma~\ref{thm:lvv_padding_prefix_mps} gives squared error at most
    $2c/N\le R^{-2}$ and bond at most
    $[P_0(N+1)^{d_0}]^2=P(N+1)^d$.
    The coefficient and degree were chosen before both $L$ and $R$.
\end{proof}
